\documentclass[11pt]{article}
\usepackage[a4paper,margin=25mm]{geometry}
\usepackage{amsmath,amssymb,amsthm,hyperref}
\newtheorem{lemma}{Lemma}\newtheorem{theorem}{Theorem}
\newcommand{\E}{\mathbb E}
\title{Global high-order polarization tails for bounded interval polynomials}
\author{Auxiliary proof; independent mathematical reviews completed}
\date{30 September 2026}
\begin{document}\maketitle
The full argument passed separate mathematical reviews by the rational-design
and independent-directions collaborators; this is not formal verification.
This note supplies the global polynomial tail estimate used in the
finite-Jackson endpoint-profile argument. It asserts only the polynomial
derivative and high-order tail estimates below, not a new profile theorem.

\begin{lemma}[All-order interior Bernstein bound]\label{bernstein}
Let $P$ be a polynomial of degree at most $d\ge1$, with
$\|P\|_{L^\infty[0,T]}\le U$. Then, for $0<v<T$ and $1\le k\le d$,
\[
 \frac{|P^{(k)}(v)|}{k!}
 \le U\left(\frac{Cd}{k\sqrt{v(T-v)}}\right)^k,           \tag{1}
\]
where $C$ is absolute.
\end{lemma}
\begin{proof}
First work on $[-1,1]$. For $z$ outside the interval put
$g(z)=\log|z+\sqrt{z^2-1}|$, choosing the root so the modulus is
at least one, and set $g=0$ on the interval. If $x\in(-1,1)$,
$a=\sqrt{1-x^2}$, and $|z-x|\le R$, then
\[
 g(z)\le C\sqrt R,\qquad g(z)\le CR/a.                  \tag{2}
\]
For the first bound when $R\le1$, compare a square root of $z^2-1$
to the closer of the two square roots of $x^2-1$. Their distance is
at most $\sqrt{|z^2-x^2|}\le C\sqrt R$, since the product of the
two candidate distances is $|z^2-x^2|$.
Both boundary values $x\pm\sqrt{x^2-1}$ have modulus one. Thus
$e^{g(z)}\le1+R+C\sqrt R$, which proves the bound. For $R\ge1$,
$g(z)\le\log(C(1+R))\le C\sqrt R$.
If $R\le a^2/8$, a local analytic branch of
$\log(z+\sqrt{z^2-1})$ has derivative of modulus
$|z^2-1|^{-1/2}\le C/a$ on the disk, because its distance from each
endpoint is comparable to the corresponding distance of $x$.
Integration from $x$ proves the second bound in this case. Otherwise
the first bound and $R>a^2/8$ prove it.

For completeness, the Bernstein--Walsh estimate here follows directly
from the maximum principle: if $Q$ is bounded by $U$ on $[-1,1]$,
then $w^{-d}Q((w+w^{-1})/2)$ is analytic for $|w|>1$, including
infinity, and bounded by $U$ on $|w|=1$. Hence
$|Q(z)|\le U e^{d g(z)}$.
Cauchy's formula on the circle of radius $R=ka/d$, using (2), gives
$|Q^{(k)}(x)|/k!\le U(Cd/(ka))^k$.
Rescaling $x=2v/T-1$ proves (1).
\end{proof}

Use the marked-tensor notation
\[
 M=m-1,\quad \ell=\lfloor M/4\rfloor,\quad N=M-\ell,
 \quad z_j=\ell(1-p_j),\quad s=\sum_{j\ne i}(1-p_j).
\]
For a uniform $\ell$-subset $B$ of the unmarked indices, let
$A_B=\prod_{j\in B}p_j$, let $R$ be its complement, and put
\[
 v=N^{-1}\sum_{j\in R}z_j,\qquad
 a_k=e_k((z_j-v)_{j\in R})/\binom Nk.
\]
We use the established endpoint balance and Laplace row bound
\[
 \max_jz_j\le C(s+1),\qquad
 \E_w e^{-c s}\le C_c/m\quad(c>0\text{ fixed}).          \tag{3}
\]
The marked column is included in the maximum. The elementary
centered-coefficient estimate gives
\[
 |a_k|\le\left(\frac{e k\sigma^2}{N^2}\right)^{k/2},
 \qquad \sigma^2=\sum_{j\in R}(z_j-v)^2.                 \tag{4}
\]

\begin{theorem}[Negligible high differential orders]\label{tail}
Fix $A,C_0>0$. There are constants $\kappa>0$ and $B>0$ such that
the following holds for all sufficiently large $m$.
Let $d=\lfloor\kappa m\rfloor$, $J=\lceil B\log m\rceil$, and
let $P$ have degree at most $d$ with
$\|P\|_{L^\infty[0,\ell]}\le m^{C_0}$.
Then
\begin{align*}
 &(\ell+1)\E_w\E_B A_B
 \sum_{k=J+1}^{d+1}|a_k|
 \left[(1+|z_i-v|)\frac{|P^{(k)}(v)|}{k!}
            +\frac{|P^{(k-1)}(v)|}{(k-1)!}\right]
 \le m^{-A}.
\end{align*}
Derivatives beyond the degree vanish. This estimate covers both the
unmarked and direct marked polarization remainders.
\end{theorem}
\begin{proof}
Write $U=\|P\|_\infty$. First restrict to subsets with $v\le\ell/2$.
From (3),
$\sigma^2\le\sum_{j\in R}z_j^2\le CN(s+1)v$.
Combining (1) and (4), for $k\ge2$, gives
\[
 |a_k|\frac{|P^{(k)}(v)|}{k!}
 \le U\left(\frac{Cd}{m}\sqrt{\frac{s+1}{k}}\right)^k.
                                                               \tag{5}
\]
Here $\ell-v\ge\ell/2\asymp m$ and $N\asymp m$.
The formulas extend to $v=0$: then all complementary $z_j$ vanish,
so $a_k=0$ for $k\ge1$.

For each row, Maclaurin's inequality gives
$\E_BA_B\le e^{-\ell s/M}$. From (3) it follows that for every
$h\ge0$ and fixed $c>0$,
\[
 (\ell+1)\E_w e^{-cs}(s+1)^h
                         \le C^{h+1}\Gamma(h+1).          \tag{6}
\]
Use (6) on (5), including the extra factor
$1+|z_i-v|\le C(s+1)$. Stirling's bounds show that the resulting
integrated $k$th term is at most
$CU(k+1)^3(Cd/m)^k$.
For the derivative of order $k-1$, the same calculation leaves one
factor $\sqrt{k(s+1)v/N}$. Since $v\le Cs$, its integrated bound is
\[
 CU(k+1)^3m^{-1/2}(Cd/m)^{k-1}.                          \tag{7}
\]
Constants can change from line to line. These estimates also cover
$k=d+1$, whose derivative of order $k$ is zero.

On the remaining subsets $v>\ell/2$, use instead the universal
variance bound
\[
 \sigma^2\le N v(\ell-v),
\]
which follows from $0\le z_j\le\ell$ and
$\sum z_j^2\le\ell\sum z_j$.
Equations (1) and (4) therefore give
\[
 |a_k|\frac{|P^{(k)}(v)|}{k!}
 \le U\left(\frac{Cd}{\sqrt{Nk}}\right)^k
 \le U\exp(Cd^2/N).                                   \tag{8}
\]
The last inequality follows by maximizing
$(k/2)\log(Cd^2/(Nk))$ over positive real $k$.
For the derivative of order $k-1$, the same argument costs at most
an additional polynomial factor in $m$, since
$\sqrt{v(\ell-v)}\le\ell/2$. At $v=\ell$ all $a_k$ for $k\ge1$
vanish, so no endpoint exception is needed.

The condition $v>\ell/2$ forces $s>N/2$, because
$v\le\ell s/N$. Hence the success-weighted contribution of these
rows is at most a polynomial factor times
\[
 U\exp(Cd^2/N-cm).
\]
Indeed one can first bound each summand by (8), then average $A_B$
using $\E_BA_B\le e^{-\ell s/M}$ and finally use $\E_w1=1$.
Choose the fixed $\kappa$ small enough to make this exponentially
small in $m$.

Decrease $\kappa$ further so the ratios in (5)--(7) after integration
are at most a fixed $\rho<1$.
The sum of those bounds over $k>J$ is at most
$CU\sum_{k>J}(k+1)^3\rho^{k-1}$.
Since $U\le m^{C_0}$, choosing $B$ sufficiently large in terms of
$A,C_0$ makes the sum at most half the required bound. The
exponentially small part is eventually smaller than the other half.
\end{proof}

This argument is entirely global over the rows and does not require
an interior lower cutoff on their deficit sum. To remove the logarithm
in the profile rate, one still needs a positive polynomial kernel of
width $m^{-1/2}$ with localized derivatives through order $J$.
\end{document}
